\section{Introducción}\label{sec:introduccion}

La corriente alterna (CA) se caracteriza por invertir periódicamente su sentido de circulación, propiedad que viabilizó la transmisión eléctrica a gran escala e impulsó, desde la Segunda Revolución Industrial, la producción en masa y el auge de las máquinas eléctricas rotativas~\cite{Chapman2012,GuiaProyecto2026}. En todo sistema de potencia en CA, el flujo energético se descompone en potencia activa ($P$), responsable del trabajo mecánico útil; potencia reactiva ($Q$), requerida para establecer los campos magnéticos de excitación; y potencia aparente ($S$), que fija el esfuerzo térmico de conductores y transformadores~\cite{Chapman2012}. La relación entre potencia útil y aparente define el factor de potencia ($\cos\phi$), indicador clave del aprovechamiento eléctrico tanto en redes monofásicas domiciliarias como en instalaciones trifásicas industriales~\cite{GuiaProyecto2026}. Un bajo $\cos\phi$, sumado a arranques ineficientes y pérdidas electromecánicas sin control, eleva la demanda sobre el Sistema Interconectado Nacional (SIN) y amplifica el impacto ambiental asociado al factor de emisión eléctrico colombiano (\qty{164.38}{\gram/(\kilo\watt\hour)} de $\text{CO}_2$)~\cite{UPME2024}.

Frente a este desafío, el proyecto responde al compromiso académico con los Objetivos de Desarrollo Sostenible ---en particular ODS~7 (energía asequible y no contaminante), ODS~9 (industria, innovación e infraestructura) y ODS~12 (producción y consumo responsables)---, en línea con la Agenda~2030 y la Misión de Sabios~\cite{GuiaProyecto2026}. Desde la perspectiva tecnológica, la propuesta materializa la transición de la Industria~4.0, centrada en la digitalización y automatización de máquinas, hacia la Industria~5.0, enfocada en el ser humano, la resiliencia operativa y la sostenibilidad ambiental~\cite{Breque2021,Xu2021}. La convergencia entre Internet Industrial de las Cosas (IIoT), gemelos digitales e Inteligencia Artificial (IA) permite convertir variables eléctricas, térmicas y vibratorias en diagnósticos oportunos que protegen simultáneamente los equipos y al personal~\cite{Singh2023,GuiaProyecto2026}.

El trabajo aborda el sector asignado de planta molino de granos (\textbf{AgroMolinos del Valle S.A.}), complejo agroindustrial de los Llanos Orientales dedicado a la recepción, trituración en molino de martillos, secado rotatorio, clasificación ciclónica, ensacado de \qty{50}{\kilogram} y despacho de cereales. Su accionamiento crítico es un motor trifásico de inducción jaula de ardilla de \qty{10}{HP} (\qty{7.46}{\kilo\watt}), \qty{460}{\volt} y \qty{60}{\hertz}, clase de eficiencia IE3~\cite{WEG2024,IEC60034}. Durante la molienda, esta máquina afronta fluctuaciones estocásticas de carga, atmósferas con polvo orgánico combustible clasificadas como ATEX Zona~22~\cite{NFPA61} y riesgo elevado de paradas imprevistas por sobrecalentamiento estatórico, desbalance de fases o desgaste en rodamientos~\cite{Moubray1997,Nandi2005}.

Para atender integralmente estas exigencias, el proyecto articula de forma progresiva los tres avances del curso en una única arquitectura ciberfísica. En el \emph{Avance~1} se efectuó la caracterización reglamentaria conforme al RETIE~\cite{RETIE2024,RETIE2026} y el modelado multifísico en MATLAB/Simulink con Control Directo de Par y lazo de velocidad PI (DTC/PI). En el \emph{Avance~2} se implementó el diagnóstico predictivo mediante un Perceptrón Multicapa (MLP $8\to 16\to 8\to 4$) y una Máquina de Vectores de Soporte (SVM RBF), integrados con una ruta de Mantenimiento Centrado en Confiabilidad (RCM)~\cite{Moubray1997} y un panel interactivo en MATLAB App Designer. En el \emph{Avance~3}, toda la solución converge en un \textbf{Juego Serio (\textit{Serious Game}) 3D interactivo en Godot Engine~4.7.2 con Vulkan~1.3 Forward+}~\cite{Singh2023,Stefan2024}, que incorpora optimización energética por variador de frecuencia frente a arranque directo (VFD vs.\ DOL), inducción física diegética de cinco fallas sin botones 2D, inducción segura de 20 estaciones con equipo de protección personal completo (EPP~6/6) y las 5 Reglas de Oro LOTO~\cite{RETIE2026}, además de una Tablet Industrial SCADA nativa de nueve pestañas sin navegador externo.

\subsection{Objetivos del Proyecto}\label{subsec:objetivos}

\textbf{Objetivo general:} Diseñar, validar e implementar una plataforma ciberfísica orientada a la Industria~5.0 para AgroMolinos del Valle S.A., integrando modelado electromecánico, diagnóstico predictivo con IA, optimización energética VFD y entrenamiento inmersivo en seguridad eléctrica bajo normativa RETIE.

\textbf{Objetivos específicos:}
\begin{enumerate}\setlength{\itemsep}{2pt}\setlength{\topsep}{2pt}
\item Modelar y validar el motor trifásico de \qty{10}{HP} bajo el marco RETIE~\cite{RETIE2024,RETIE2026}, articulando la simulación multifísica DTC/PI en MATLAB/Simulink con la ecuación de Kloss en tiempo real.
\item Entrenar y desplegar clasificadores supervisados (MLP $8\to 16\to 8\to 4$ y SVM RBF) con exactitud superior al \qty{98}{\percent}, integrados a una ruta técnica de mantenimiento RCM.
\item Optimizar el consumo energético mediante control VFD frente a DOL, superando un ahorro del \qty{25}{\percent} en energía específica (\unit{\kilo\watt\hour/\tonne}) con mitigación cuantificable de $\text{CO}_2\text{e}$~\cite{UPME2024}.
\item Implementar y certificar, mediante QA automatizado 10/10, el Serious Game 3D en Godot~4.7.2 con inducción segura de 20 estaciones, protocolo LOTO, cinco fallas diegéticas y Tablet SCADA nativa~\cite{Stefan2024}.
\end{enumerate}
